% Mathematical Addendum — Response to Prof. Fierro's critique (Q1-Q3)
% Canonical source. The version as prepared 2026-06-29 is archived as
% addendum_received_2026-06-29.pdf. Amendments F1/F2/F4a/F4c from
% REVIEW_2026-07-11.md APPLIED (ratified by rmw3 2026-07-11).
% FIERRO-NS-RANGE: N and sigma-range FILLED (EXP-NS-01 primary data).
% N=3000 ticks per noise level; sigma max=4.0||x|| (σ_in=2.0, noise=4×norm).
% FIERRO-L-BOUND: computed — see REVIEW_2026-07-11.md (L_f > 1 for all τ).
\documentclass[11pt]{article}
\usepackage[margin=1.1in]{geometry}
\usepackage{amsmath, amssymb, amsthm}
\usepackage[T1]{fontenc}

\title{Mathematical Addendum: Response to Critique on the Feasible\\
Region $\Omega$ and Stochastic Dynamics}
\author{VSF Research Group}
\date{July 2026 (rev.\ 2)}

\newcommand{\R}{\mathbb{R}}
\newcommand{\E}{\mathbb{E}}

\begin{document}
\maketitle

\section*{Context}

We thank the reviewer for identifying three structural gaps in the original
existence and invariance argument for the feasible region
$\Omega \subset \R^{128}$ under stochastic updates
$x(t+1) = f_0(x(t)) + \sigma\eta(t)$. All three points have been formally
resolved and integrated into the revised manuscript (Sec.~3--4). Below is the
consolidated mathematical statement.

\section*{1. Convexity of $\Omega$}

The global region $\Omega = \bigcap_i \{x : A_i(x)\}$ is not convex: the
entropy bound $H(x) \le H_{\max}$ is an upper level set of a concave function
(non-convex in general), and the ethical exclusion $A_0$ is a non-geometric
predicate (covenant-hash validity) that need not preserve convexity. A direct
application of Brouwer's fixed-point theorem on $\Omega$ is therefore
unjustified.

\paragraph{Resolution.} The reviewer is right --- and the corrected argument
no longer requires convexity at any step. It has three independent legs:
\begin{enumerate}
  \item \textbf{Existence of an invariant measure (Krylov--Bogolyubov).}
        The state space is a compact set and the update kernel is Feller
        ($f_0$ continuous, additive noise), so an invariant probability
        measure exists (Kryloff \& Bogoliouboff [1937]). Compactness
        suffices; convexity plays no role.
  \item \textbf{Deterministic fixed point (Banach).} On the contractive core
        \[
          \Omega_A = \bigl\{x \in \Omega : \langle x, s\rangle \ge
          \theta_A \lVert x\rVert\,\lVert s\rVert\bigr\},
          \qquad \theta_A = 0.88,
        \]
        --- the intersection of a circular cone (convex for $\theta_A \ge 0$)
        with the bounded state set, nested inside the viable region
        $\Omega \supseteq \{\cos \ge 0.70\}$ --- the measured contraction
        $\lambda_{\max} = \ln(0.99) < 0$ (EXP-LY-01), an average per-step
        log-contraction along core trajectories, makes some iterate
        $f_0^{\,n}$ a contraction; the Banach principle applied to
        $f_0^{\,n}$ then gives a unique fixed point of $f_0$. The argument
        requires completeness and invariance of the core, but not convexity.
  \item \textbf{Uniqueness and rate (Foster--Lyapunov).} Uniqueness of the
        invariant measure and geometric convergence follow from the drift
        condition of Sec.~3 below (Meyn \& Tweedie [2009]).
\end{enumerate}
The sum-form route via Krasnosel'ski\u{\i}'s theorem for a contraction plus a
compact continuous perturbation (Krasnosel'ski\u{\i} [1955]) remains
available, but it is no longer load-bearing: no step of the corrected
argument depends on the convexity of $\Omega$ or of $\Omega_A$.

\section*{2. Continuity and Invariance of $f_0$}

Forcing $f_0(\Omega) \subseteq \Omega$ by projection would distort the natural
dynamics and obscure failure modes. We do not impose invariance globally.

\paragraph{Resolution.} The deterministic map $f_0$ exhibits layered
behavior, and we are explicit about what is measured and what is not claimed:
\begin{itemize}
  \item On $\Omega_A$: the measured $\lambda_{\max} = \ln(0.99) < 0$
        certifies average per-step log-contraction along core trajectories
        (equivalently, some iterate of $f_0$ is a contraction --- the form
        used in Sec.~1). The global worst-case Lipschitz bound computable
        from the deployed constants,
        $L_f \le \max_s \lVert\nabla \ell_s\rVert \cdot
        \operatorname{diam}\{\mu_s\}/\tau$, is not informative on the core:
        there the softmax mixture is near-saturated and its local gradient
        lies far below the worst case.
  \item On $\Omega \setminus \Omega_A$: we do \emph{not} claim per-step
        non-expansiveness. The same worst-case bound evaluates to
        $L_f \in [1.4,\ 24.4]$ across the deployed temperature range, so a
        per-step bound $L \le 1$ cannot be discharged from the code
        constants. Confinement off the core is carried empirically instead:
        excursions recover within a measured finite time
        $\tau_{\mathrm{rec}}$ (sub-30\,s across 20 forced thermal cycles),
        and sustained noise produces zero exits from $\Omega$ (Sec.~3).
        Beyond the zero-exit observation, the deployed kernel's viability
        margin has been \emph{measured directly}: over $15{,}435$ samples
        across $10.7$\,h of continuous runtime the distance to the viable
        boundary is $0.097$ (median), strictly positive and stable
        (max $0.100$), with viability holding on $100\%$ of samples.
        Confinement off the core is therefore quantified by a positive
        measured margin, not merely inferred from the absence of exits.
\end{itemize}
No artificial clamping is applied. The geometry of the Lyapunov function
$L(x)$ ensures that states perturbed outside $\Omega_A$ are naturally
re-contracted into the core without leaving $\Omega$ (Remark~2, Sec.~3).
Invariance is thus a \emph{measured dynamical property}, not an imposed
algebraic constraint.

\section*{3. Stochastic Perturbations $\sigma \cdot \eta(t)$}

Even if $f_0 : \Omega \to \Omega$ were invariant, the additive noise
$\sigma\eta(t)$ with unbounded support breaks pointwise invariance:
$x(t+1) \notin \Omega$ with non-zero probability.

\paragraph{Resolution.} Two observations. First, in the deployed system the
perturbation is bounded by construction: the state lives on a compact set and
the noise amplitude is governed by the $A_{\mathrm{breath}}$ axiom, so the
unbounded-support premise applies to the idealized model, not to the running
kernel. Second, for the idealized model we replace pointwise invariance with
\emph{stochastic viability} (Aubin [1991]) and the Lyapunov drift condition
\[
  \E\bigl[L(x(t+1)) - L(x(t)) \mid x(t) \in \Omega\bigr] \le 0,
\]
which is exactly the geometric-ergodicity condition for the induced Markov
operator (Meyn \& Tweedie [2009]). Confinement is then measured: EXP-NS-01
drives the kernel with sustained bounded noise at every tick and observes
\emph{zero exit events} over $N = 3000$ ticks per noise level
(empirical exit rate $< 3/N$ at 95\% confidence)
for all $\sigma \in [0,\, 4.0\,\|x\|]$.
The viability margin $m(\sigma)$ degrades smoothly and saturates at
$\approx 0.33$; no finite noise ceiling was observed. The system does not rely
on hard boundaries; it relies on negative expected drift that recontracts
excursions in finite time.
% FILLED 2026-07-11 (FIERRO-NS-RANGE): N=3000, σ_max=4.0||x||.
% EXP-NS-01 primary data: sweep σ_in∈[0,2.0], label "noise=4×distribution norm".
% The addendum claimed 2.0||x|| → corrected to 4.0||x|| (matches manuscript Remark).

\section*{Integration with the Manuscript}

The three objections are resolved in the current revision:
\begin{enumerate}
  \item \textbf{Sec.~3, Def.~2 \& Thm.~1:} Existence carried by
        Krylov--Bogolyubov (compactness + Feller) and the Banach principle on
        the core; convexity no longer enters the argument at any point.
  \item \textbf{Sec.~4, Thm.~2 \& EXP-LY-01:} Convergence in expectation via
        measured Lyapunov drift ($\lambda_{\max} = -0.010050$), not assumed
        invariance.
  \item \textbf{Sec.~6.4, EXP-NS-01:} Sustained-noise margin curve replaces
        the borrowed impulsive proxy; zero observed exits up to
        $\sigma = 4.0\,\|x\|$ verify stochastic viability empirically.
\end{enumerate}

The existence proof is now \emph{measured, not postulated}: contraction,
confinement, and ergodicity are discharged by hardware experiments rather than
asserted algebraically. Brouwer's theorem is no longer invoked; the negative
Lyapunov exponent and the viability-theoretic reading of $\Omega$ carry the
argument.

\begin{thebibliography}{9}
\bibitem{aubin} J.-P. Aubin, \emph{Viability Theory}, Birkh\"auser, 1991.
\bibitem{kras} M.~A. Krasnosel'ski\u{\i}, ``Two remarks on the method of
successive approximations,'' \emph{Uspekhi Mat.\ Nauk}, 10(1):123--127, 1955.
\bibitem{krylov} N.~Kryloff and N.~Bogoliouboff, ``La th\'eorie g\'en\'erale
de la mesure dans son application \`a l'\'etude des syst\`emes dynamiques de
la m\'ecanique non lin\'eaire,'' \emph{Ann.\ of Math.}, 38(1):65--113, 1937.
\bibitem{meyn} S.~P. Meyn and R.~L. Tweedie, \emph{Markov Chains and
Stochastic Stability}, Cambridge Univ.\ Press, 2nd ed., 2009.
\end{thebibliography}

\end{document}
